\documentclass[UTF8]{ctexart}
\usepackage{geometry}
\usepackage{amsmath,amssymb,amsfonts}
\usepackage{array,booktabs}
\geometry{a4paper,margin=2.5cm}

\begin{document}

\begin{center}
\LARGE\bfseries 概统作业题单
\end{center}

\section*{教材《概率论及其应用》——1.8习题（p.20/21）}

\begin{enumerate}
\renewcommand{\labelenumi}{\textbf{\arabic{enumi}.}}
\setcounter{enumi}{13}
\item 试验证下列关系式\footnote{注意：$(A\cup B)'$表示$A\cup B$的补集，与$A'\cup B'$的意义不同。同样，$(AB)'$与$A'B'$的意义亦不相同。}：
\begin{enumerate}
\item[(a)] $(A\cup B)' = A'B'$;
\item[(b)] $(A\cup B)-B = A-AB = AB'$;
\item[(c)] $AA = A\cup A = A$;
\item[(d)] $(A-AB)\cup B = A\cup B$;
\item[(e)] $(A\cup B)-AB = AB'\cup A'B$;
\item[(f)] $A'\cup B' = (AB)'$;
\item[(g)] $(A\cup B)C = AC\cup BC$.
\end{enumerate}

\setcounter{enumi}{15}
\item 试述下列关系中哪些是正确的，哪些是错误的：
\begin{enumerate}
\item[(a)] $(A\cup B)-C = A\cup(B-C)$;
\item[(b)] $ABC = AB(C\cup B)$;
\item[(c)] $A\cup B\cup C = A\cup(B-AB)\cup(C-AC)$;
\item[(d)] $A\cup B = (A-AB)\cup B$;
\item[(e)] $(AB\cup BC\cup CA) \supset ABC$;
\item[(f)] $(AB\cup BC\cup CA) \subset (A\cup B\cup C)$;
\item[(g)] $(A\cup B)-A = B$;
\item[(h)] $AB'C \subset A\cup B$;
\item[(i)] $(A\cup B\cup C)' = A'B'C'$;
\item[(j)] $(A\cup B)'C = A'C\cup B'C$;
\item[(k)] $(A\cup B)'C = A'B'C$;
\item[(l)] $(A\cup B)'C = C-C(A\cup B)$.
\end{enumerate}

\item 令$A,B,C$为任意三个事件。试用$A,B,C$来表达下列事件：
\begin{enumerate}
\item[(a)] 只有$A$发生；
\item[(b)] $A$和$B$都发生而$C$不发生；
\item[(c)] 所有这三个事件都发生；
\item[(d)] $A,B,C$中至少有一个事件发生；
\item[(e)] 至少有两个事件发生；
\item[(f)] 恰有一个事件发生；
\item[(g)] 恰有两个事件发生；
\item[(h)] 没有一个事件发生；
\item[(i)] 不多于两个事件发生。
\end{enumerate}

\setcounter{enumi}{18}
\item 利用习题18的结果证明
\[
P\{A\cup B\cup C\} = P\{A\}+P\{B\}+P\{C\}-P\{AB\}-P\{AC\}-P\{BC\}+P\{ABC\}.
\]
{[这是(4.1.5)的特殊情形。]}
\end{enumerate}

\section*{教材《概率论及其应用》——2.10习题（pp.45--47）}

\begin{enumerate}
\renewcommand{\labelenumi}{\textbf{\arabic{enumi}.}}
\setcounter{enumi}{8}
\item $n$个球随机放入$n$个盒中，问恰有一盒空着的概率为多少？

\setcounter{enumi}{10}
\item 某人有$n$把钥匙，其中只有一把能打开他的门。他逐个地用它们去试开（抽样是无放回的），这可能要$1,2,\cdots,n$次才能把门打开。证明这$n$种不同结果的概率都是$1/n$。

\setcounter{enumi}{20}
\item 谣言的传播。在拥有$n+1$个居民的城市里，某人告诉第二个人一个谣言，第二个人又把谣言告诉第三个人，如此等等。在每一步中，谣言的接收者都是随机地从$n$个居民中挑选的。求下述两事件的概率：谣言传播了$r$次后，(a)还没有回到第一个造谣者；(b)没有一个人两次听到谣言。当每一次都把谣言同时告诉由城市中随机选取的$N$个居民时，问上面两事件的概率等于多少（前面的问题是$N=1$的特别情形）。

\setcounter{enumi}{23}
\item 求下述两事件的概率：(a)12个人的生日在12个不同的月份（假定任何一个人生于12个月中之任一月都是等概的）；(b)6个人的生日恰巧在两个月中。

\setcounter{enumi}{37}
\item 印刷错误。假设书中每一页都有$N$个印刷符号，它们有可能印错。全书共有$n=500$页和$r=50$个印错的符号。证明：(a)第$1,2,\cdots,n$页分别含有$r_1,r_2,\cdots,r_n$个印错的符号的概率等于
\[
\binom{N}{r_1}\binom{N}{r_2}\cdots\binom{N}{r_n}\Big/\binom{nN}{r}.
\]
(b)当$N$充分大，这个概率可以用(2.5.3)来逼近，并推出$r$个错误分布在$n$页的问题近似于$r$个球随机地放入$n$个盒的问题。（注意：这可作为费米-狄拉克统计的一个普遍极限性质，见2.5.1节。）
\end{enumerate}

\section*{教材《概率论及其应用》——5.8习题（pp.116--118）}

\begin{enumerate}
\renewcommand{\labelenumi}{\textbf{\arabic{enumi}.}}
\setcounter{enumi}{5}
\item 厂里有三台机器$A,B,C$生产螺丝钉，它们的产量分别占总产量的25\%、35\%、40\%，而产品的次品率分别为5\%、4\%、2\%。今从产品中任意取出一颗螺丝钉，发现它是次品。求这颗螺丝钉是由机器$A$（或$B$或$C$）造出的概率。

\setcounter{enumi}{6}
\item 设100个男人中有5个色盲，而1000个女人中有25个色盲。今从人群中任选一人，发现他（她）是色盲。求此人是男性的概率（假设人群中男女的人数是相同的）。

\setcounter{enumi}{17}
\item 波利亚罐子模型，5.2节例(c)。若已知第2次所取出的球是黑的，求第1次所取出的球也是黑的概率。

\setcounter{enumi}{18}
\item 波利亚罐子模型，5.2节例(c)。用数学归纳法证明：无论是哪一次试验，黑球出现的概率总是$b/(b+r)$。

\setcounter{enumi}{25}
\item 两两独立但非相互独立。掷两颗骰子并定义如下三个事件：$A$表示第一颗骰子掷出奇数点，$B$表示第二颗骰子掷出奇数点，$C$表示两颗骰子的点数之和为奇数（即一颗掷出偶数，另一颗掷出奇数）。假定36个样本点中每个的概率是$1/36$，则任何两个事件皆相互独立，且每个概率都是$1/2$。但这三个事件不能同时发生。
\end{enumerate}

\section*{教材《概率论及其应用》——6.10习题（p.139）}

\begin{enumerate}
\renewcommand{\labelenumi}{\textbf{\arabic{enumi}.}}
\setcounter{enumi}{4}
\item 设中靶的概率是$\dfrac{1}{5}$，问10次独立的打靶中至少有2次中靶的概率是多少？

\setcounter{enumi}{11}
\item 如果总体中色盲的比例是1\%，问有放回的随机样本应该多大才能使样本中有1个色盲的概率大于等于0.95？

\setcounter{enumi}{15}
\item 一本书包含$n$页，每一页的印错处平均为$\lambda$。计算至少有1页其印错之处超过$k$的概率。
\end{enumerate}

\section*{教材《概率论与数理统计》（盛聚等）——第一章习题（pp.25--29）}

\begin{enumerate}
\renewcommand{\labelenumi}{\textbf{\arabic{enumi}.}}
\setcounter{enumi}{1}
\item 设$A,B,C$为三个事件，用$A,B,C$的运算关系表示下列各事件：
\begin{enumerate}
\item[(1)] $A$发生，$B$与$C$不发生。
\item[(2)] $A$与$B$都发生，而$C$不发生。
\item[(3)] $A,B,C$中至少有一个发生。
\item[(4)] $A,B,C$都发生。
\item[(5)] $A,B,C$都不发生。
\item[(6)] $A,B,C$中不多于一个发生。
\item[(7)] $A,B,C$中不多于两个发生。
\item[(8)] $A,B,C$中至少有两个发生。
\end{enumerate}

\setcounter{enumi}{7}
\item 在1500件产品中有400件次品、1100件正品。任取200件。
\begin{enumerate}
\item[(1)] 求恰有90件次品的概率。
\item[(2)] 求至少有2件次品的概率。
\end{enumerate}

\setcounter{enumi}{11}
\item 50只铆钉随机地取来用在10个部件上，其中有3只铆钉强度太弱。每个部件用3只铆钉。若将3只强度太弱的铆钉都装在一个部件上，则这个部件强度就太弱。问发生一个部件强度太弱的概率是多少？

\item 一俱乐部有5名一年级学生，2名二年级学生，3名三年级学生，2名四年级学生。
\begin{enumerate}
\item[(1)] 在其中任选4名学生，求一、二、三、四年级的学生各一名的概率。
\item[(2)] 在其中任选5名学生，求一、二、三、四年级的学生均包含在内的概率。
\end{enumerate}

\setcounter{enumi}{15}
\item 据以往资料表明，某一3口之家，患某种传染病的概率有以下规律：
\[
P\{\text{孩子得病}\}=0.6,\quad P\{\text{母亲得病}\mid\text{孩子得病}\}=0.5,
\]
\[
P\{\text{父亲得病}\mid\text{母亲及孩子得病}\}=0.4,
\]
求母亲及孩子得病但父亲未得病的概率。

\setcounter{enumi}{17}
\item 某人忘记了电话号码的最后一个数字，因而他随意地拨号。求他拨号不超过三次而接通所需电话的概率。若已知最后一个数字是奇数，则此概率是多少？

\item 
\begin{enumerate}
\item[(1)] 设甲袋中装有$n$只白球、$m$只红球，乙袋中装有$N$只白球、$M$只红球。今从甲袋中任意取一只球放入乙袋中，再从乙袋中任意取一只球。问取到白球的概率是多少？
\item[(2)] 第一只盒子装有5只红球、4只白球，第二只盒子装有4只红球、5只白球。先从第一盒中任取2只球放入第二盒中去，然后从第二盒中任取一只球。求取到白球的概率。
\end{enumerate}

\setcounter{enumi}{22}
\item 将两信息分别编码为A和B传送出去，接收站收到时，A被误收作B的概率为0.02，而B被误收作A的概率为0.01。信息A与信息B传送的频繁程度为$2:1$。若接收站收到的信息是A，问原发信息是A的概率是多少？

\setcounter{enumi}{24}
\item 某人下午5:00下班，他所积累的资料表明：
\begin{center}
\begin{tabular}{|c|c|c|c|c|c|}
\hline
到家时间 & 5:35--5:39 & 5:40--5:44 & 5:45--5:49 & 5:50--5:54 & 迟于5:54 \\
\hline
乘地铁的概率 & 0.10 & 0.25 & 0.45 & 0.15 & 0.05 \\
\hline
乘汽车的概率 & 0.30 & 0.35 & 0.20 & 0.10 & 0.05 \\
\hline
\end{tabular}
\end{center}
某日他抛一枚硬币决定乘地铁还是乘汽车，结果他是5:47到家的。试求他乘地铁回家的概率。

\item 病树的主人外出，委托邻居浇水，设已知如果不浇水，树死去的概率为0.8。若浇水则树死去的概率为0.15。有0.9的把握确定邻居会记得浇水。
\begin{enumerate}
\item[(1)] 求主人回来时树还活着的概率。
\item[(2)] 若主人回来时树已死去，求邻居忘记浇水的概率。
\end{enumerate}

\setcounter{enumi}{28}
\item 根据报道美国人血型的分布近似地为：A型为37\%，O型为44\%，B型为13\%，AB型为6\%。夫妻拥有的血型是相互独立的。
\begin{enumerate}
\item[(1)] B型的人只有输入B,O两种血型才安全。若妻为B型，夫为何种血型未知，求夫是妻的安全输血者的概率。
\item[(2)] 随机地取一对夫妇，求妻为B型，夫为A型的概率。
\item[(3)] 随机地取一对夫妇，求其中一人为A型，另一人为B型的概率。
\item[(4)] 随机地取一对夫妇，求其中至少有一人是O型的概率。
\end{enumerate}

\setcounter{enumi}{32}
\item 盒中有编号为1,2,3,4的4只球，随机地自盒中取一只球，事件$A$为``取得的是1号或2号球''，事件$B$为``取得的是1号或3号球''，事件$C$为``取得的是1号或4号球''。验证：
\[
P(AB)=P(A)P(B),\quad P(AC)=P(A)P(C),\quad P(BC)=P(B)P(C),
\]
但
\[
P(ABC)\ne P(A)P(B)P(C),
\]
即事件$A,B,C$两两独立，但$A,B,C$不是相互独立的。

\item 如果一危险情况C发生时，一电路闭合并发出警报，我们可以借用两个或多个开关并联以改善可靠性。在C发生时这些开关每一个都应闭合，且若至少一个开关闭合了，警报就发出。如果两个这样的开关并联连接，它们每个具有0.96的可靠性（即在情况C发生时闭合的概率），问这时系统的可靠性（即电路闭合的概率）是多少？如果需要有一个可靠性至少为0.9999的系统，则至少需要用多少只开关并联？设各开关闭合与否是相互独立的。

\setcounter{enumi}{34}
\item 设$A,B,C$是三个事件，且$P(A)=P(B)=P(C)=\dfrac{1}{4}$，$P(AB)=P(BC)=0$，$P(AC)=\dfrac{1}{8}$。求$A,B,C$中至少有一个发生的概率。

\setcounter{enumi}{37}
\item 袋中装有$m$枚正品硬币、$n$枚次品硬币（次品硬币的两面均印有国徽），在袋中任取一枚，将它投掷$r$次，已知每次都得到国徽。问这枚硬币是正品的概率为多少？
\end{enumerate}

\end{document}
